% Standalone problem document, generated from the adjacent README.md.
% Compile with XeLaTeX. Mathematical content is maintained in Markdown.
\documentclass[11pt,a4paper]{article}
\usepackage[fontsize=10.5pt]{scrextend}
\usepackage[margin=22mm,headheight=15pt,headsep=9mm,footskip=11mm]{geometry}
\usepackage{fontspec}
\setmainfont{texgyrepagella}[Extension=.otf,UprightFont=*-regular,BoldFont=*-bold,ItalicFont=*-italic,BoldItalicFont=*-bolditalic]
\setsansfont{texgyreheros}[Extension=.otf,UprightFont=*-regular,BoldFont=*-bold,ItalicFont=*-italic,BoldItalicFont=*-bolditalic]
\setmonofont{texgyrecursor}[Extension=.otf,UprightFont=*-regular,BoldFont=*-bold,ItalicFont=*-italic,BoldItalicFont=*-bolditalic,Scale=MatchLowercase]
\usepackage{amsmath,amssymb}
\usepackage{unicode-math}
\setmathfont{texgyrepagella-math.otf}
\usepackage{xcolor,microtype,enumitem,needspace,fancyhdr}
\usepackage{bookmark}
\definecolor{ink}{HTML}{17344B}
\definecolor{accent}{HTML}{197D80}
\definecolor{muted}{HTML}{596773}
\hypersetup{colorlinks=true,linkcolor=accent,urlcolor=accent,pdftitle={MI-31 - Sharp
parameter dependence for structured Gaussian operator
norms},pdfauthor={Open Problems in Numerical Linear Algebra}}
\urlstyle{same}
\setlength{\parindent}{0pt}
\setlength{\parskip}{5pt plus 1pt minus 1pt}
\linespread{1.04}
\setlength{\emergencystretch}{3em}
\widowpenalty=10000
\clubpenalty=10000
\displaywidowpenalty=10000
\allowdisplaybreaks[1]
\setlist{nosep,leftmargin=1.5em}
\providecommand{\tightlist}{\setlength{\itemsep}{0pt}\setlength{\parskip}{0pt}}
\providecommand{\pandocbounded}[1]{#1}
\pagestyle{fancy}
\fancyhf{}
\fancyhead[L]{\sffamily\footnotesize\color{muted}OPEN PROBLEMS IN NUMERICAL LINEAR ALGEBRA}
\fancyhead[R]{\sffamily\bfseries\footnotesize\color{accent}MI-31}
\fancyfoot[L]{\parbox[b]{0.9\textwidth}{\sffamily\fontsize{8}{10}\selectfont\color{muted}Editorial ratings; a bounded search cannot certify that a problem remains unsolved.\\Literature check: 2026-09-11}}
\fancyfoot[R]{\sffamily\footnotesize\color{muted}\thepage}
\renewcommand{\headrulewidth}{0.3pt}
\renewcommand{\headrule}{\hbox to\headwidth{\color{accent}\leaders\hrule height \headrulewidth\hfill}}
\makeatletter
\renewcommand{\section}{\Needspace{5\baselineskip}\@startsection{section}{1}{0pt}{12pt plus 2pt minus 2pt}{4pt}{\sffamily\bfseries\large\color{ink}}}
\renewcommand{\subsection}{\Needspace{4\baselineskip}\@startsection{subsection}{2}{0pt}{9pt plus 1pt minus 1pt}{3pt}{\sffamily\bfseries\normalsize\color{ink}}}
\makeatother
\setcounter{secnumdepth}{0}
\begin{document}
{\sffamily\small\color{muted}Matrix inequalities and norms\par}
\vspace{3pt}
{\raggedright\hyphenpenalty=10000\exhyphenpenalty=10000\sffamily\bfseries\fontsize{21}{24}\selectfont\color{ink}Sharp
parameter dependence for structured Gaussian operator norms\par}
\vspace{7pt}
{\sffamily\small\textbf{Difficulty:} challenging\quad\textbf{Importance:} interesting
to the community\par}
{\sffamily\small\bfseries\color{accent}Status: Partially resolved\par}
\vspace{4pt}
{\color{accent}\hrule height 0.8pt}
\vspace{6pt}
\textbf{Topic:} Expected norms of random matrices with unequal entry
variances

\textbf{Rating rationale:} The expectation is characterized for fixed
norm exponents, but making the constant uniform requires sharper control
as the exponents and dimensions vary. Such estimates quantify
amplification by random matrices under different input and output norms.

\subsection{Statement}\label{statement}

Let \(m,n\) be positive integers, \(1\le p\le2\le q\le\infty\), and
\(A=(a_{ij})\in\mathbb R^{m\times n}\). Let \(p^*\) satisfy
\(1/p+1/p^*=1\), with \(p^*=\infty\) when \(p=1\), and set
\(L(k)=\max\{1,\ln k\}\) for integers \(k\ge1\). For independent
standard real Gaussian variables \(g_{ij}\), set \(G_A=(a_{ij}g_{ij})\).
Define

\[
\|B\|_{p\to q}=\sup_{\|x\|_p\le1}\|Bx\|_q,
\qquad
D_1=\max_{1\le i\le m}\|(a_{ij})_{j=1}^n\|_{p^*},
\qquad
D_2=\max_{1\le j\le n}\|(a_{ij})_{i=1}^m\|_q.
\]

\textbf{Conjecture (Latała--Strzelecka).} There is an absolute constant
\(C>0\), independent of \(m,n,p,q,A\), such that

\[
\mathbb E\|G_A\|_{p\to q}
\le C\left[
\sqrt{\min\{p^*,L(n)\}}\,D_1
+\sqrt{\min\{q,L(m)\}}\,D_2
+\mathbb E\max_{i,j}|a_{ij}g_{ij}|
\right].
\]

Use the usual maximum norm for an \(\ell_\infty\) vector norm and
\(\min\{\infty,L(k)\}=L(k)\). The conjecture is a uniform upper bound;
it does not assert a matching lower bound with these same three terms
for every variance profile.

\subsection{Known cases and numerical
significance}\label{known-cases-and-numerical-significance}

The spectral case \(p=q=2\) is known. For \(a_{ij}=1\) throughout, the
two dimension-capped square-root factors have the correct order. The
same paper proves comparison with
\(D_1+D_2+\mathbb E\max_{i,j}|a_{ij}g_{ij}|\) using constants allowed to
depend on \(p,q\). That theorem resolves the older
Guédon--Hinrichs--Litvak--Prochno conjecture, while the displayed
universal constant remains the new target.

The matrix \(G_A\) models independent random entry errors with a
prescribed variance profile. Its induced norm measures worst-case
amplification from an input \(\ell_p\) norm to an output \(\ell_q\)
norm. The parameter dependence matters when the norms themselves vary
with dimension.

\newpage
\subsection{References and status
check}\label{references-and-status-check}

\begin{itemize}
\tightlist
\item
  R. Latała and M. Strzelecka, \emph{Operator \(\ell_p\to\ell_q\) norms
  of Gaussian matrices}, Advances in Mathematics 501 (2026), article
  111097. \href{https://doi.org/10.1016/j.aim.2026.111097}{DOI};
  \href{https://arxiv.org/html/2502.02186v3}{current arXiv v3}, dated
  2026-06-09. Conjecture 5 is the displayed upper bound; the notation
  preceding Conjecture 1 defines the absolute implicit constant and the
  truncated logarithm. Theorem 2 and Remarks 3--4 distinguish the proved
  comparison from the conjectured parameter dependence.
\item
  R. Latała, R. van Handel and P. Youssef, \emph{The dimension-free
  structure of nonhomogeneous random matrices}, Inventiones Mathematicae
  214 (2018), 1031--1080. The spectral-norm result is identified as
  reference 15 in the preceding source.
\end{itemize}

On 2026-09-11, checked the current arXiv version and published full
text, the authors' institutional publication records, and exact-title,
author/Conjecture-5, sharp-parameter-dependence, proof and
counterexample searches through 2026. No full resolution of Conjecture 5
was located. Statements that the older Conjecture 1 has been solved do
not settle this uniform-constant refinement. This is a bounded
literature check.
\end{document}
